\section{基于KG-GCN的燃煤机组汽轮机系统故障诊断方法}
\label{sec:kg_gcn_v4}

\subsection{引言}

汽轮机热力系统故障往往同时作用于多个相互关联的状态参数。对于同一故障，其影响可沿主通流、再热、抽汽回热及冷端边界向相邻设备传播；对于不同故障，又可能由于共享局部参数而表现出相似的残差变化。因此，仅根据单个参数幅值或局部时序特征进行诊断，难以充分利用设备之间客观存在的热力联系。第三章从机组设备结构、DCS测点资料以及公开故障知识中构建了汽轮机系统故障知识图谱，为不同设备和参数之间的关联关系提供了结构化表达。

第二章建立的健康状态模型能够根据机组当前运行边界给出设备状态参数的正常参考值。将实际DCS测量值与健康模型输出之间的偏差作为动态状态信息，并进一步映射到知识图谱对应参数节点，可以同时利用运行数据和静态领域知识描述设备故障。基于此，本文在课题组既有KG-GCN研究思路\cite{YangHaodong2026CFB}基础上，采用图卷积网络对知识图谱约束下的多参数残差进行逐层聚合，并利用无标签残差窗口实施掩码预训练，在少量带标签故障样本条件下完成故障分类，从而构建面向燃煤机组汽轮机系统的KG-GCN故障诊断方法。

\subsection{图卷积网络}
\label{subsec:gcn_theory_v4}

复杂汽轮机热力系统包含主通流、再热、抽汽以及冷端等多种连接关系，其参数之间并非规则网格上的相邻关系，而是由设备结构和介质传递路径决定的非欧几里得拓扑。传统神经网络通常直接在规则向量或矩阵空间内提取特征，难以显式表达不同参数之间的结构依赖。图卷积网络（Graph Convolutional Network，GCN）能够在图拓扑上对节点自身及其邻域信息进行聚合，使节点特征在关联路径上进行传递，从而适用于具有明确设备连接关系的工业系统\cite{Kipf2017GCN}。

设参数图为$G=(V,E)$，其中$V=\{v_1,v_2,\ldots,v_N\}$为参数节点集合，$E$为节点之间的关联边。图的拓扑连接关系由邻接矩阵$A\in\mathbb{R}^{N\times N}$表示，节点初始特征矩阵记为$H^{(0)}$。为在图卷积过程中同时保留节点自身信息，在邻接矩阵中加入单位矩阵$I_N$，得到
\begin{equation}
\tilde A=A+I_N.
\end{equation}
相应度矩阵为
\begin{equation}
\tilde D_{ii}=\sum_j\tilde A_{ij}.
\end{equation}
对邻接矩阵进行对称归一化，得到
\begin{equation}
\bar A=\tilde D^{-\frac{1}{2}}\tilde A\tilde D^{-\frac{1}{2}}.
\end{equation}
标准图卷积层的特征传播形式可写为
\begin{equation}
H^{(l+1)}=\sigma\left(\bar A H^{(l)}W^{(l)}\right),
\end{equation}
式中：$H^{(l)}$为第$l$层节点特征；$W^{(l)}$为该层可学习权重矩阵；$\sigma(\cdot)$为非线性激活函数。由式可知，每一层图卷积均依据邻接矩阵确定的信息传播路径，将目标节点自身与相邻节点的状态共同映射至新的特征空间。随着图卷积层数增加，节点能够逐步获得更远邻域内的关联信息，从而形成对系统局部故障传播特征的高阶表示。

本文采用两层GCN完成参数关联信息聚合。考虑汽轮机固定DCS节点自身残差同样具有重要诊断意义，在图卷积计算中保留节点自身表示，并采用残差连接抑制多层传播导致的节点信息衰减。其计算形式为
\begin{equation}
H^{(1)}=\mathrm{ReLU}\left[\mathrm{LN}\left(H^{(0)}+\bar A H^{(0)}W^{(1)}\right)\right],
\end{equation}
\begin{equation}
H^{(2)}=\mathrm{ReLU}\left[\mathrm{LN}\left(H^{(1)}+\bar A H^{(1)}W^{(2)}\right)\right],
\end{equation}
式中：$\mathrm{LN}(\cdot)$表示LayerNorm操作。第一层GCN主要融合直接相邻参数的信息，第二层进一步聚合两跳邻域特征，使节点表示同时包含自身残差变化和局部热力拓扑信息。

\subsection{KG-GCN模型输入构建}
\label{subsec:kg_input_v4}

第三章构建的汽轮机系统故障知识图谱同时包含设备、参数、异常、故障及处置等多种类型实体。为将知识图谱用于图卷积故障诊断，需要首先从系统级知识图谱中提取能够与实际DCS状态参数对应的参数节点，并依据设备归属、介质流向及直接热力接口关系形成参数关联图。设由系统知识图谱投影得到的诊断参数图为
\begin{equation}
G_p=(V_p,E_p),
\end{equation}
式中：$V_p$为可与DCS状态参数建立映射的节点集合；$E_p$为由知识图谱结构关系投影得到的参数关联边。若节点$v_i$与$v_j$之间存在允许用于状态传播的知识关系，则参数邻接矩阵$A_p$中对应元素取1，否则取0。

对于第$i$个状态参数，在时刻$t$的实际测量值和健康状态模型输出分别记为$y_i(t)$和$\hat y_i(t)$。定义二者之间的状态偏差为
\begin{equation}
e_i(t)=y_i(t)-\hat y_i(t).
\end{equation}
在设备正常运行时，健康模型能够解释由负荷、主蒸汽边界和操作调节引起的大部分正常状态变化，因此$e_i(t)$主要体现模型误差和测量扰动；当设备出现异常后，实际运行状态偏离健康参考状态，相关参数残差将产生持续或协同变化。由此，状态监测残差可作为后续故障诊断的动态状态输入。

由于不同参数量纲和幅值范围存在差异，采用训练集健康残差统计量进行标准化：
\begin{equation}
r_i(t)=\frac{e_i(t)-\mu_i}{\sigma_i+\varepsilon},
\end{equation}
式中：$\mu_i$和$\sigma_i$分别为训练集第$i$个参数健康残差的均值和标准差；$\varepsilon$为防止分母为零的微小常数。验证集和测试集不参与标准化参数估计。

故障演化具有连续时间特征，因此对连续$L_r$个采样时刻构造残差窗口。时刻$t$对应的输入矩阵可表示为
\begin{equation}
R_t=
\begin{bmatrix}
r_1(t-L_r+1)&\cdots&r_1(t)\\
r_2(t-L_r+1)&\cdots&r_2(t)\\
\vdots&\ddots&\vdots\\
r_N(t-L_r+1)&\cdots&r_N(t)
\end{bmatrix}
\in\mathbb{R}^{N\times L_r}.
\end{equation}
矩阵的每一行对应知识图谱中的一个固定参数节点，每一列对应连续时间位置。由此，原本独立排列的多参数残差被转换为具有明确图节点身份的动态特征矩阵。第五章VHP代表性算例中，最终诊断参数节点数$N=24$，残差窗口长度取$L_r=60$。

在输入层中，首先将每个参数节点的残差窗口映射至$d$维初始特征空间：
\begin{equation}
H^{(0)}=\mathrm{ReLU}(R_tW_e+b_e),
\end{equation}
其中$W_e\in\mathbb{R}^{L_r\times d}$。随后，将$H^{(0)}$与由知识图谱得到的归一化邻接矩阵$\bar A_p$共同输入GCN，通过多层邻域聚合获得具有知识结构约束的节点特征。对于本文固定DCS测点，节点所在的物理位置本身具有诊断意义，因此在形成图级特征时保留节点固定顺序，并将节点局部表示与图传播表示组合后展开为图级特征向量$\mathbf g$，用于后续分类。

\subsection{KG-GCN模型训练}
\label{subsec:kg_training_v4}

实际电站运行过程中，正常和一般变工况数据能够长期连续积累，而高质量故障样本数量相对有限。若直接依赖少量带标签故障样本训练图卷积网络，模型容易过度拟合有限故障模式。为此，本文采用无标签残差窗口掩码预训练与有标签故障分类相结合的两阶段训练方式，使图卷积编码器首先学习知识图谱结构下不同参数残差之间的关联规律，再利用带标签故障窗口训练分类层。

设无标签残差窗口样本集合为
\begin{equation}
\mathcal D_u=\{R^{(1)},R^{(2)},\ldots,R^{(M_u)}\},
\end{equation}
其中$M_u$为无标签残差窗口数量。预训练阶段对输入残差窗口随机遮蔽。设掩码矩阵为$M\in\{0,1\}^{N\times L_r}$，其中$M_{ij}=0$表示第$i$个参数在第$j$个时间位置的残差被遮蔽，$M_{ij}=1$表示该位置保持可见，则遮蔽后的残差窗口为
\begin{equation}
\tilde R=M\odot R+(1-M)\odot R_{mask},
\end{equation}
式中：$\odot$表示逐元素乘法；$R_{mask}$为用于替代被遮蔽位置的掩码值，本文取0。实验中随机遮蔽比例为15\%。

将$\tilde R$与归一化邻接矩阵$\bar A_p$输入图卷积编码器，得到
\begin{equation}
Z=E_{\theta}(\bar A_p,\tilde R),
\end{equation}
式中：$E_{\theta}$为待预训练的GCN编码器；$\theta$表示其中可学习参数。随后利用解码器将结构化特征还原至原始残差窗口空间：
\begin{equation}
\hat R=D_{\phi}(Z),
\end{equation}
式中：$D_{\phi}$为解码器。预训练阶段仅在被遮蔽位置计算重构误差，其损失函数为
\begin{equation}
\mathcal L_{mask}=\frac{\sum_{ij}(1-M_{ij})(R_{ij}-\hat R_{ij})^2}{\sum_{ij}(1-M_{ij})+\varepsilon}.
\end{equation}
通过最小化掩码重构损失，GCN编码器需要利用未遮蔽节点状态和知识图谱邻域关系恢复缺失信息，从而学习参数残差在图结构中的关联特征。

完成无标签残差窗口预训练后，固定图卷积编码器参数，并利用有标签故障窗口训练分类层。设有标签故障样本集合为
\begin{equation}
\mathcal D_l=\{(R^{(n)},y^{(n)})\}_{n=1}^{M_l},
\end{equation}
式中：$R^{(n)}$为第$n$个故障残差窗口；$y^{(n)}$为对应故障类别标签；$M_l$为带标签故障样本数量。将故障残差窗口输入预训练后的图卷积编码器，得到图级特征
\begin{equation}
\mathbf g^{(n)}=F\left(E_{\theta^*}(\bar A_p,R^{(n)})\right),
\end{equation}
式中：$\theta^*$表示完成掩码预训练后的固定编码器参数；$F(\cdot)$表示图级特征整理操作。

分类层采用全连接映射与Softmax函数输出各故障类别概率：
\begin{equation}
\hat{\mathbf p}^{(n)}=\mathrm{Softmax}(W_c\mathbf g^{(n)}+b_c).
\end{equation}
若故障类别数为$C$，则$\hat{\mathbf p}^{(n)}$中的每一个元素表示该残差窗口属于对应类别的预测概率。分类阶段采用交叉熵损失函数：
\begin{equation}
\mathcal L_{cls}=-\frac{1}{M_l}\sum_{n=1}^{M_l}\sum_{k=1}^{C}y_{nk}\log \hat p_{nk}.
\end{equation}
通过上述训练过程，大量无标签状态残差首先用于学习参数间结构关联，随后少量带标签故障样本用于确定故障类别判别边界。完成训练后，KG-GCN即可对新的状态残差窗口输出故障类别，并结合知识图谱中的参数连接关系对诊断结果进行解释。

\subsection{本章小结}

本章在第三章汽轮机系统故障知识图谱基础上，构建了基于KG-GCN的故障诊断模型。首先介绍GCN在图拓扑上的特征传播机制，并利用汽轮机系统知识关系形成可与DCS参数对应的邻接矩阵；随后定义实际测量值与健康状态模型输出之间的残差，构造参数--时间残差窗口并映射到知识图谱参数节点；最后采用随机掩码重构完成无标签预训练，在固定图卷积编码器的基础上利用带标签故障窗口训练Softmax分类层。由此形成“状态监测残差--知识图谱节点映射--图卷积特征聚合--掩码预训练--故障分类”的诊断过程。下一章选取VHP作为代表性设备开展算例验证。
